% IEEE Transactions on Intelligent Transportation Systems LaTeX Template
% 用法: 直接复制此文件，替换 [YOUR CONTENT] 部分

\documentclass[journal,10pt]{IEEEtran}

% ===== 常用包 =====
\usepackage[utf8]{inputenc}
\usepackage[T1]{fontenc}
\usepackage{amsmath,amssymb,amsfonts}
\usepackage{graphicx}
\usepackage{cite}
\usepackage{algorithmic}
\usepackage{algorithm}
\usepackage{booktabs}
\usepackage{multirow}
\usepackage{color}
\usepackage{url}
\usepackage{hyperref}
\usepackage[caption=false,font=footnotesize]{subfig}

% ===== 图表路径 =====
\graphicspath{{figures/}}

% ===== 标题和作者 =====
\title{[YOUR PAPER TITLE]}

\author{
\IEEEauthorblockN{Author One\IEEEauthorrefmark{1},
Author Two\IEEEauthorrefmark{1},
Author Three\IEEEauthorrefmark{2}}
\IEEEauthorblockA{\IEEEauthorrefmark{1}School of Transportation, University Name, City, Country}
\IEEEauthorblockA{\IEEEauthorrefmark{2}Department of Computer Science, University Name, City, Country}
\thanks{Corresponding author: Author One (email@university.edu)}
}

% ===== 开始文档 =====
\begin{document}

\maketitle

% ===== Abstract =====
\begin{abstract}
[YOUR ABSTRACT HERE - 150-200 words]

Traffic flow prediction is a critical task for intelligent transportation systems.
However, existing methods often fail to capture [LIMITATION].
In this paper, we propose [METHOD] to address this challenge.
[1-2 sentences describing the approach].
Extensive experiments on [DATASETS] demonstrate that our method
achieves [RESULTS]. [1-2 sentences about broader impact].
\end{abstract}

\begin{IEEEkeywords}
Traffic flow prediction, spatiotemporal graph neural network,
intelligent transportation systems.
\end{IEEEkeywords}

% ===== I. Introduction =====
\section{Introduction}
\label{sec:intro}

[YOUR INTRODUCTION - 6段漏斗结构]

% 第1段：领域重要性
Traffic flow prediction is a critical task for intelligent
transportation systems (ITS), enabling efficient traffic management
and route planning~\cite{ref1}.

% 第2段：技术挑战
However, accurate traffic prediction remains challenging due to
[CHALLENGES]~\cite{ref2}.

% 第3段：现有方法分类
Existing methods can be categorized into [CATEGORIES]~\cite{ref3,ref4}.

% 第4段：现有方法的局限
Despite significant progress, existing methods suffer from [LIMITATIONS].

% 第5段：本文方案
To address these challenges, we propose [METHOD]. Our key insight is [INSIGHT].

% 第6段：贡献列表
The main contributions of this paper are summarized as follows:
\begin{itemize}
    \item We propose [CONTRIBUTION 1].
    \item We design [CONTRIBUTION 2].
    \item Extensive experiments on [DATASETS] demonstrate [CONTRIBUTION 3].
\end{itemize}

% ===== II. Related Work =====
\section{Related Work}
\label{sec:related}

\subsection{Graph Neural Networks for Traffic Prediction}
[PARAGRAPH ABOUT GNN METHODS]

\subsection{Transformer-based Methods}
[PARAGRAPH ABOUT TRANSFORMER METHODS]

\subsection{Other Approaches}
[PARAGRAPH ABOUT OTHER METHODS]

% ===== III. Method =====
\section{Proposed Method}
\label{sec:method}

\subsection{Problem Formulation}
\label{sec:problem}

Given a traffic network represented as a graph $\mathcal{G} = (\mathcal{V}, \mathcal{E}, \mathbf{A})$,
where $\mathcal{V}$ is the set of $N$ sensors, $\mathcal{E}$ is the set of edges,
and $\mathbf{A} \in \mathbb{R}^{N \times N}$ is the adjacency matrix.
Let $\mathbf{X}^{(t)} \in \mathbb{R}^{N \times F}$ be the traffic observations at time $t$,
where $F$ is the number of features.
The goal is to predict future traffic states $\hat{\mathbf{X}}^{(t+1:t+T)}$
given historical observations $\mathbf{X}^{(t-L+1:t)}$,
where $L$ is the input length and $T$ is the prediction horizon.

\subsection{Overall Architecture}
\label{sec:architecture}

The proposed model consists of [N] main components:
(a) [COMPONENT 1], (b) [COMPONENT 2], and (c) [COMPONENT 3].

% 架构图
\begin{figure}[!t]
    \centering
    \includegraphics[width=\columnwidth]{figures/architecture.pdf}
    \caption{Overall architecture of the proposed [MODEL NAME].
    The model consists of (a) [Component 1], (b) [Component 2],
    and (c) [Component 3].}
    \label{fig:architecture}
\end{figure}

\subsection{Module 1: [MODULE NAME]}
\label{sec:module1}

[MODULE DESCRIPTION WITH MATH]

\begin{equation}
\mathbf{Y} = \text{Attention}(\mathbf{Q}, \mathbf{K}, \mathbf{V})
= \text{softmax}\left(\frac{\mathbf{Q}\mathbf{K}^T}{\sqrt{d_k}}\right)\mathbf{V}
\label{eq:attention}
\end{equation}

where $\mathbf{Q}$, $\mathbf{K}$, and $\mathbf{V}$ are query, key, and value
matrices, respectively, and $d_k$ is the dimension of the keys.

\subsection{Module 2: [MODULE NAME]}
\label{sec:module2}

[MODULE DESCRIPTION]

\subsection{Training Objective}
\label{sec:loss}

The model is trained to minimize the mean absolute error (MAE):
\begin{equation}
\mathcal{L} = \frac{1}{N \cdot T} \sum_{i=1}^{N} \sum_{t=1}^{T}
\left| \hat{x}_i^{(t)} - x_i^{(t)} \right|
\label{eq:loss}
\end{equation}

% ===== IV. Experiments =====
\section{Experiments}
\label{sec:experiments}

\subsection{Experimental Setup}
\label{sec:setup}

\subsubsection{Datasets}
We evaluate our method on four widely-used traffic benchmark datasets:
\begin{itemize}
    \item \textbf{METR-LA}: 207 sensors on Los Angeles highways,
    traffic speed, 5-minute intervals.
    \item \textbf{PEMS-BAY}: 325 sensors in San Francisco Bay Area,
    traffic speed, 5-minute intervals.
    \item \textbf{PEMS04}: 307 sensors in San Bernardino,
    traffic flow, 5-minute intervals.
    \item \textbf{PEMS08}: 170 sensors in San Bernardino,
    traffic flow, 5-minute intervals.
\end{itemize}

\subsubsection{Baselines}
We compare our method with the following baselines:
\begin{itemize}
    \item \textbf{DCRNN}~\cite{dcrnn}: Diffusion Convolutional Recurrent Neural Network.
    \item \textbf{STGCN}~\cite{stgcn}: Spatio-Temporal Graph Convolutional Network.
    \item \textbf{Graph WaveNet}~\cite{gwnet}: Graph WaveNet with adaptive adjacency matrix.
    \item \textbf{STAEformer}~\cite{staeformer}: Spatio-Temporal Adaptive Embedding Transformer.
    \item \textbf{PDFormer}~\cite{pdformer}: Propagation Delay-Aware Dynamic Long-Range Transformer.
\end{itemize}

\subsubsection{Implementation Details}
\label{sec:impl}

Our model is implemented in PyTorch and trained on a single NVIDIA A100 GPU.
We use the Adam optimizer with an initial learning rate of $1 \times 10^{-3}$
and a batch size of 64. The input sequence length is 12 (1 hour),
and the prediction horizons are 15, 30, and 60 minutes.
We use Z-score normalization with statistics computed on the training set.
All experiments are repeated 5 times with different random seeds,
and we report the mean and standard deviation.

\subsection{Main Results}
\label{sec:results}

Table~\ref{tab:main} shows the comparison results on [YOUR DATASET 1] and [YOUR DATASET 2] datasets.
Our method achieves the best performance across all metrics and prediction horizons.

% WARNING: Replace ALL placeholder values below with your actual experimental results.
% DO NOT use these example numbers in your paper.

\begin{table*}[!t]
\centering
\caption{Performance comparison on [YOUR DATASET 1] and [YOUR DATASET 2] datasets.
The best results are in \textbf{bold} and the second best are \underline{underlined}.
↓ indicates lower is better.}
\label{tab:main}
\begin{tabular}{l|ccc|ccc}
\toprule
\multirow{2}{*}{Method} & \multicolumn{3}{c|}{[YOUR DATASET 1]} & \multicolumn{3}{c}{[YOUR DATASET 2]} \\
\cmidrule{2-7}
& MAE↓ 15min & MAE↓ 30min & MAE↓ 60min & MAE↓ 15min & MAE↓ 30min & MAE↓ 60min \\
\midrule
[Baseline 1] & [YOUR DATA] & [YOUR DATA] & [YOUR DATA] & [YOUR DATA] & [YOUR DATA] & [YOUR DATA] \\
[Baseline 2] & [YOUR DATA] & [YOUR DATA] & [YOUR DATA] & [YOUR DATA] & [YOUR DATA] & [YOUR DATA] \\
[Baseline 3] & [YOUR DATA] & [YOUR DATA] & [YOUR DATA] & [YOUR DATA] & [YOUR DATA] & [YOUR DATA] \\
\midrule
\textbf{Ours} & \textbf{[YOUR DATA]} & \textbf{[YOUR DATA]} & \textbf{[YOUR DATA]} & \textbf{[YOUR DATA]} & \textbf{[YOUR DATA]} & \textbf{[YOUR DATA]} \\
\bottomrule
\end{tabular}
\end{table*}

\subsection{Ablation Study}
\label{sec:ablation}

To validate the effectiveness of each component, we conduct ablation studies
on METR-LA with 60-minute prediction horizon.

% WARNING: Replace ALL placeholder values below with your actual experimental results.

\begin{table}[!t]
\centering
\caption{Ablation study on [YOUR DATASET] ([YOUR HORIZON]).
w/o means removing the component.}
\label{tab:ablation}
\begin{tabular}{lccc}
\toprule
Variant & MAE↓ & RMSE↓ & MAPE↓ \\
\midrule
Full Model & \textbf{[YOUR DATA]} & \textbf{[YOUR DATA]} & \textbf{[YOUR DATA]} \\
w/o Spatial Attention & 3.12 & 6.15 & 8.01\% \\
w/o Temporal Attention & 3.18 & 6.28 & 8.25\% \\
w/o Adaptive Embedding & 3.05 & 5.98 & 7.78\% \\
\bottomrule
\end{tabular}
\end{table}

The results show that removing the temporal attention module leads to the
largest performance degradation (+0.20 MAE), indicating its importance
in capturing long-range temporal dependencies.

\subsection{Efficiency Analysis}
\label{sec:efficiency}

Table~\ref{tab:efficiency} compares the computational efficiency of different methods.

\begin{table}[!t]
\centering
\caption{Computational efficiency comparison.}
\label{tab:efficiency}
\begin{tabular}{lcccc}
\toprule
Method & Params (M) & FLOPs (G) & Time (ms) & MAE \\
\midrule
STGCN & 0.5 & 2.1 & 3.2 & 4.54 \\
GWNet & 0.8 & 5.3 & 8.5 & 3.53 \\
STAEformer & 8.2 & 15.6 & 25.3 & 3.15 \\
\midrule
Ours & 5.1 & 10.2 & 15.8 & 2.98 \\
\bottomrule
\end{tabular}
\end{table}

\subsection{Visualization}
\label{sec:visualization}

Figure~\ref{fig:prediction} shows the prediction results on two representative sensors.

\begin{figure}[!t]
    \centering
    \includegraphics[width=\columnwidth]{figures/prediction.pdf}
    \caption{Prediction results on METR-LA.
    (a) Sensor with stable traffic patterns.
    (b) Sensor with high variability.
    Black: ground truth. Red: our prediction.}
    \label{fig:prediction}
\end{figure}

% ===== V. Conclusion =====
\section{Conclusion}
\label{sec:conclusion}

In this paper, we proposed [METHOD] for traffic flow prediction.
[1-2 sentences summarizing the approach].
Extensive experiments on four benchmark datasets demonstrated that
our method achieves state-of-the-art performance, with improvements
of [X]\% in MAE over the strongest baseline.
[1-2 sentences about limitations and future work].

% ===== Acknowledgments =====
\section*{Acknowledgments}
This work was supported by [FUNDING].
The authors thank [ACKNOWLEDGMENTS].

% ===== References =====
\bibliographystyle{IEEEtran}
\bibliography{references}

\end{document}
